% Drawing 57/08-5 -- folder 57, batch 08, archive page 154 (PDF page 155).
% Transcription: en marge gauche, une suite de croquis de courbes : une boucle à point double, des configurations de droites concourantes, un point marqué « 1 » ; d'autres croquis en bas de page
% Kind: margin sketches of curves. Method 2 (schematic redraw from the facsimile). Status: ok.
% Notes: Left margin, top to bottom. A loop with a double point: the long stroke (doubled at its top) runs down through the heavy double point, the loop leaves it and closes back on it, and two short strokes cross the loop on the right. Three crossing lines. Two lines through two heavy points: one line through both points, one through the upper point (ending in a hook), one through the lower point; the mark "1" at the upper right. Every stroke through a heavy point is drawn through the same computed point. The small dot inside the loop is left out (a speck). The "other sketches at the bottom of the page" are those of 08-6 and are drawn there. Redrawn 2026-11-09 (issue 128): to check again.
% Coordinates: 1 unit = 1 mm = 1/100 of the crop width (bbox in page fractions: [0.1, 0.24, 0.28, 0.53]).
% Redrawn by Claude (claude-opus-5-5), issue #73, 2026-11-09; not checked by a human; not compiled (no TeX in the session).
% Rights: reproduces part of the work -- behind RENDER_TRANSCRIPTS, never published by itself.
\documentclass[tikz,border=3pt]{standalone}
\usepackage{amsmath,amssymb}
\usetikzlibrary{patterns,arrows.meta}
\tikzset{pencil/.style={black!45}, struck/.style={black!35}, mark/.style={circle,draw,inner sep=1.3pt}, pt/.style={circle,fill,inner sep=0.9pt}, opt/.style={circle,draw,fill=white,inner sep=1.1pt}}
\begin{document}
\begin{tikzpicture}[x=1mm,y=1mm,line cap=round,line join=round,font=\scriptsize]
\draw plot[smooth, tension=0.55] coordinates {(19.4,-11.4) (21.1,-28.5) (26.5,-42.5) (22.8,-54.8) (24.5,-72.1)};
\draw plot[smooth, tension=0.55] coordinates {(20,-14) (20.5,-22) (23,-33)};
\draw plot[smooth, tension=0.55] coordinates {(26.5,-42.5) (30.3,-30.8) (36.5,-20.3) (44,-22.8) (47.9,-37.7) (45.7,-51.4) (39.4,-55.7) (32.5,-52.5) (26.5,-42.5)};
\node[pt] at (26.5,-42.5) {};
\draw (42.8,-34) -- (56.5,-29.1);
\draw (40,-44) -- (53.9,-48.9);
\draw (57.1,-83.3) -- (45.4,-111.3);
\draw (42.8,-103.3) -- (65.1,-129);
\draw plot[smooth, tension=0.55] coordinates {(72.5,-101) (54.8,-115.9) (37.7,-137) (20.5,-167.2)};
\draw plot[smooth, tension=0.55] coordinates {(63.4,-169.5) (71.3,-188.4) (78.8,-204.3) (86.2,-214.6)};
\draw plot[smooth, tension=0.55] coordinates {(85,-177) (71.3,-188.4) (47,-213.5) (38.2,-228.2)};
\draw plot[smooth, tension=0.55] coordinates {(26.8,-201.5) (47,-213.5) (62.8,-228.2)};
\node[pt] at (71.3,-188.4) {};
\node[pt] at (47,-213.5) {};
\node[inner sep=1pt] at (87.3,-167.2) {\small $1$};
\end{tikzpicture}
\end{document}
